\section{The LAMP Protocol}
\label{sec:construction}

This section gives the formal construction of \(\protocol\), following the
overview in Section~\ref{sec:tech-overview}.  The protocol proves a matrix
multiplication statement over committed encoded data.  The SNARK circuit checks
only sampled codeword positions, while Merkle openings and CP-Link proofs bind
the sampled circuit witnesses to values fixed before the query positions are
derived.


\subsection{Committed Statement}
\label{sec:relation-chain}

Let \(\mat{A},\mat{B},\mat{C}\in\mathbb{F}^{k\times k}\), and let
\(\mathcal{C}:\mathbb{F}^k\to\mathbb{F}^n\) be a systematic linear code with
relative distance \(\delta\).  The target relation is
\[
  \mathcal{R}_{\mathsf{MM}}
  =
  \{(\mat{A},\mat{B},\mat{C}) :
    \mat{A}\mat{B}=\mat{C}\}.
\]
The verifier does not receive the raw matrices as public input.  Instead, the
prover commits to row-wise encodings of the matrices and proves that the
committed objects satisfy the multiplication relation.

As in Section~\ref{sec:prelim-freivalds}, \(\protocol\) derives one scalar
\(r\in\mathbb{F}\) and uses the structured Freivalds vector
\[
  \vect{r}=(1,r,r^2,\ldots,r^{k-1}).
\]
The prover then defines
\[
  \vect{x}=\vect{r}\mat{A},
  \qquad
  \vect{y}=\vect{x}\mat{B},
  \qquad
  \vect{z}=\vect{r}\mat{C}.
\]
If \(\mat{A}\mat{B}=\mat{C}\), then \(\vect{y}=\vect{z}\).  Therefore, the
matrix multiplication statement can be checked by verifying the three
vector-product relations above and the final equality \(\vect{y}=\vect{z}\).

\subsection{Matrix Encoding and Commitment}
\label{sec:matrix-commitment}

The prover first encodes the input and output matrices row-wise:
\[
  \widehat{\mat{A}}=\enc(\mat{A}),\qquad
  \widehat{\mat{B}}=\enc(\mat{B}),\qquad
  \widehat{\mat{C}}=\enc(\mat{C}).
\]
For each codeword position \(j\in[n]\), it packs the three encoded columns into
one vector
\[
\begin{aligned}
  \vect{m}_{ABC,j}
  \coloneqq
  \big(
    \widehat{\mat{A}}[*,j]\,
    \|\,\widehat{\mat{B}}[*,j]\,
    \|\,\widehat{\mat{C}}[*,j]
  \big)
  \in\mathbb{F}^{3k}.
\end{aligned}
\]
Using a Pedersen commitment key \(\ck_{ABC}\) for vectors in
\(\mathbb{F}^{3k}\), the prover samples \(o_{ABC,j}\sample\mathbb{F}\) and
computes
\[
  \cm_{ABC,j}
  \gets
  \pedC(\ck_{ABC},\vect{m}_{ABC,j};o_{ABC,j}).
\]
The matrix root is the Merkle commitment to the vector of Pedersen
commitments:
\[
  \rt_{ABC}
  \gets
  \merkleC((\cm_{ABC,j})_{j\in[n]};\bot).
\]
The root \(\rt_{ABC}\) is fixed before the Freivalds challenge \(r\) is
derived.  This ordering fixes the encoded matrix columns before the prover sees
the challenge used to form \(\vect{x},\vect{y},\vect{z}\).

\subsection{Intermediate Vectors and Query Sampling}
\label{sec:intermediate-vectors-and-query-sampling}

After receiving the matrix commitment root \(\rt_{ABC}\), the verifier samples
the Freivalds challenge \(r\sample\mathbb{F}\) and sends it to the prover.
The prover then sets $\vect{r}=(1,r,\ldots,r^{k-1})$
and computes
\[
  \vect{x}=\vect{r}\mat{A},
  \qquad
  \vect{y}=\vect{x}\mat{B},
  \qquad
  \vect{z}=\vect{r}\mat{C}.
\]
It then encodes the intermediate vectors:
\[
  \widehat{\vect{x}}=\enc(\vect{x}),
  \qquad
  \widehat{\vect{y}}=\enc(\vect{y}),
  \qquad
  \widehat{\vect{z}}=\enc(\vect{z}).
\]
For each \(j\in[n]\), the prover packs the encoded vector symbols into
\[
  \vect{m}_{XYZ,j}
  \coloneqq
  \big(
    \widehat{\vect{x}}[j]\,
    \|\,\widehat{\vect{y}}[j]\,
    \|\,\widehat{\vect{z}}[j]
  \big)
  \in\mathbb{F}^{3}.
\]
Using a Pedersen commitment key \(\ck_{XYZ}\), the prover samples
\(o_{XYZ,j}\sample\mathbb{F}\) and computes
\[
  \cm_{XYZ,j}
  \gets
  \pedC(\ck_{XYZ},\vect{m}_{XYZ,j};o_{XYZ,j}),
\]
and sends the Merkle root
\[
  \rt_{XYZ}
  \gets
  \merkleC((\cm_{XYZ,j})_{j\in[n]};\bot)
\]
to the verifier.

Only after receiving \(\rt_{XYZ}\) does the verifier sample the query tuple
\[
  I=(j_1,\ldots,j_t)\sample [n]^t
\]
and send \(I\) to the prover. In addition to checking the sampled encoded
symbols, the protocol must ensure that the encoded vector witnesses used in the
circuit are indeed encodings of the corresponding intermediate vectors. We
enforce this by barycentric out-of-domain consistency checks. To this end, the
verifier samples
\[
  \alpha_x,\alpha_y,\alpha_z \sample \mathbb{F}\setminus D,
\]
where \(D\) denotes the code evaluation domain, and sends them to the prover.
These points are later used in the code-consistency checks
\((\mathrm{CODE\text{-}X})\), \((\mathrm{CODE\text{-}Y})\), and
\((\mathrm{CODE\text{-}Z})\) of the SNARK relation. All these challenges are
public coins and are sampled only after the corresponding commitments have been
fixed.




\subsection{SNARK Relation}
\label{sec:snark-relation}
The inner SNARK proves the sampled arithmetic relation. Its public statement contains the commitment roots, the verifier-sampled challenges \(r,I,\alpha_x,\alpha_y,\alpha_z\), and the CP-SNARK witness commitments:
\[
\begin{aligned}
  \stmt =
  \big(
  &\rt_{ABC},\rt_{XYZ},r,I,\alpha_x,\alpha_y,\alpha_z,\\
  &\widetilde{\cm}_{ABC},\widetilde{\cm}_{XYZ}
  \big).
\end{aligned}
\]
The private witness contains the intermediate vectors, their codewords, and the
sampled packed values:
\[
\begin{aligned}
  \wit =
  \big(
  &\vect{x},\vect{y},\vect{z},
  \widehat{\vect{x}},\widehat{\vect{y}},\widehat{\vect{z}},\\
  &\{\vect{m}_{ABC,j}\}_{j\in I},
  \{\vect{m}_{XYZ,j}\}_{j\in I}
  \big).
\end{aligned}
\]
The CP-SNARK backend commits to the sampled witnesses
\(\{\vect{m}_{ABC,j}\}_{j\in I}\) and \(\{\vect{m}_{XYZ,j}\}_{j\in I}\), producing \(\widetilde{\cm}_{ABC}\) and \(\widetilde{\cm}_{XYZ}\).
These commitments are later linked to the external Pedersen leaves.

For every sampled \(j\in I\), write
\[
  \vect{m}_{ABC,j}
  =
  (\vect{a}_j\,\|\,\vect{b}_j\,\|\,\vect{c}_j),
  \qquad
  \vect{m}_{XYZ,j}
  =
  (x_j\,\|\,y_j\,\|\,z_j),
\]
where \(\vect{a}_j,\vect{b}_j,\vect{c}_j\in\mathbb{F}^k\) and
\(x_j,y_j,z_j\in\mathbb{F}\).  The SNARK checks:
\begingroup
\scriptsize
\begin{align}
  &\widehat{\vect{x}}
    \text{ is consistent with } \enc(\vect{x})
    \text{ at } \alpha_x,
    \tag{CODE-X}\label{eq:rs-check}\\
  &\widehat{\vect{y}}
    \text{ is consistent with } \enc(\vect{y})
    \text{ at } \alpha_y,
    \tag{CODE-Y}\label{eq:rs-y-check}\\
  &\widehat{\vect{z}}
    \text{ is consistent with } \enc(\vect{z})
    \text{ at } \alpha_z,
    \tag{CODE-Z}\label{eq:rs-z-check}\\
  &x_j=\widehat{\vect{x}}[j],
    \qquad
    y_j=\widehat{\vect{y}}[j],
    \qquad
    z_j=\widehat{\vect{z}}[j],
    \tag{LOOKUP}\label{eq:lookup-check}\\
  &x_j=\fold(\vect{r},\vect{a}_j),
    \tag{FOLD-A}\label{eq:fold-A}\\
  &y_j=\fold(\vect{x},\vect{b}_j),
    \tag{FOLD-B}\label{eq:fold-B}\\
  &z_j=\fold(\vect{r},\vect{c}_j),
    \tag{FOLD-C}\label{eq:fold-C}\\
  &y_j=z_j.
    \tag{EQ}\label{eq:eq-check}
\end{align}
\endgroup
The first three checks enforce that the intermediate codewords used by the
sampled relation are encodings of the vector witnesses.  The remaining checks
are the sampled fold relations from Section~\ref{sec:overview-encoding}.

\subsection{Merkle Openings and CP-Link}
\label{sec:openings-and-cplink}

The SNARK relation alone does not prove that the sampled witnesses are the
values committed in the Merkle trees.  The auxiliary verifier therefore checks
Merkle openings and CP-Link proofs.

For each \(j\in I\), the prover opens the external leaves
\(\cm_{ABC,j}\) and \(\cm_{XYZ,j}\) under \(\rt_{ABC}\) and \(\rt_{XYZ}\):
\begingroup
\scriptsize
\begin{align}
  &\merkleV(\rt_{ABC},j,\cm_{ABC,j},\pi_{ABC,j})=1,
    \tag{PATH-ABC}\label{eq:path-check}\\
  &\merkleV(\rt_{XYZ},j,\cm_{XYZ,j},\pi_{XYZ,j})=1.
    \tag{PATH-XYZ}\label{eq:path-xyz-check}
\end{align}
\endgroup
It then proves, via CP-Link, that the SNARK-side commitment to the sampled
\(ABC\) witness opens to the same block messages as the external Pedersen
commitments \((\cm_{ABC,j})_{j\in I}\), and similarly for the \(XYZ\) sampled
witnesses:
\begingroup
\scriptsize
\begin{align}
  &\bPi.\verify(
    \widetilde{\cm}_{ABC},
    (\cm_{ABC,j})_{j\in I},
    \pi_{\mathsf{link}}^{ABC})=1,
    \tag{LINK-ABC}\label{eq:link-check}\\
  &\bPi.\verify(
    \widetilde{\cm}_{XYZ},
    (\cm_{XYZ,j})_{j\in I},
    \pi_{\mathsf{link}}^{XYZ})=1.
    \tag{LINK-XYZ}\label{eq:link-xyz-check}
\end{align}
\endgroup
% Here \(\cpLinkV\) denotes the batched, blockwise linking relation supplied by
% the backend: the SNARK-side commitment is opened as the concatenation of the
% sampled blocks, and each external Pedersen commitment is opened to the
% corresponding block.  Equivalently, it is a batched application of the
% same-message commitment-linking relation from Section~\ref{sec:prelim-cplink}.

\begin{definition}[Sampled-opening backend]\label{def:codecom-backend}
A sampled-opening backend for \(\protocol\) consists of the Merkle opening
checks and CP-Link checks above, together with the committed-witness interface
of the SNARK.  It must satisfy:
\begin{itemize}
  \item \emph{Position binding}: an accepted Merkle opening fixes a unique
  Pedersen leaf at the queried position under the root;
  \item \emph{Commitment binding}: an accepted Pedersen leaf is binding to a
  unique message, except with negligible probability;
  \item \emph{SNARK-side binding}: an accepting SNARK proof binds the sampled
  witnesses used by the circuit to the public SNARK-side commitments
  \(\widetilde{\cm}_{ABC}\) and \(\widetilde{\cm}_{XYZ}\);
  \item \emph{CP-Link consistency}: accepted CP-Link proofs imply that the
  SNARK-side commitments and the external Pedersen leaves commit to the same
  sampled messages, except with the CP-Link soundness error.
\end{itemize}
\end{definition}

Under this backend, the verifier obtains the consistency chain
\[
\begin{array}{c}
\text{values used in the SNARK}
\\
\Updownarrow\ \text{committed witness and CP-Link}
\\
\text{external Pedersen leaves}
\\
\Updownarrow\ \text{Merkle authentication}
\\
\rt_{ABC},\rt_{XYZ}.
\end{array}
\]
This is why the SNARK circuit can avoid in-circuit Pedersen verification and
in-circuit Merkle path verification.

\subsection{Protocol Algorithms and Complexity}
\label{sec:protocol-desc}
\label{sec:complexity}

Figure~\ref{fig:protocol} summarizes the non-interactive protocol obtained by
Fiat--Shamir.  The interactive version is obtained by having the verifier
sample \(r\) after \(\rt_{ABC}\) and \(I\) after \(\rt_{XYZ}\).

\begin{figure*}[ht!]
\centering
\scalebox{1}{%
    \fbox{%
    \procedure[]{\protocol}{%
    \mathcal{P}(\mat{A},\mat{B},\mat{C}) \< \< \mathcal{V} \\
    \widehat{\mat{A}},\widehat{\mat{B}},\widehat{\mat{C}} \gets \enc(\mat{A}),\enc(\mat{B}),\enc(\mat{C}) \in \mathbb{F}^{k \times n} \< \< \\
    \text{for each } j\in[n]: \\ 
    \quad \vect{m}_{ABC,j} \coloneqq \big( \widehat{\mat{A}}[*,j]\, \|\,\widehat{\mat{B}}[*,j]\, \|\,\widehat{\mat{C}}[*,j] \big) \in\mathbb{F}^{3k} \\
    \quad o_{ABC,j}\sample\mathbb{F} \\
    \quad \cm_{ABC,j} \gets \pedC(\ck_{ABC},\vect{m}_{ABC,j};o_{ABC,j}) \\
    \rt_{ABC} \gets \merkleC((\cm_{ABC,j})_{j\in[n]};\bot) \\
    \< \sendmessageright*[2.5cm]{\rt_{ABC}} \< \\
    \< \< r \sample \FF  \\
    \< \sendmessageleft*[2.5cm]{r} \<  \\
    \vect{r}\coloneqq(1,r,\ldots,r^{k-1}) \\ \vect{x}\gets\vect{r}\mat{A}, \qquad \vect{y}\gets\vect{x}\mat{B}, \qquad \vect{z}\gets\vect{r}\mat{C} \\ \widehat{\vect{x}},\widehat{\vect{y}},\widehat{\vect{z}} \gets \enc(\vect{x}),\enc(\vect{y}),\enc(\vect{z}) \\ \text{for each } j\in[n]: \\ \quad \vect{m}_{XYZ,j} \coloneqq \big( \widehat{\vect{x}}[j]\, \|\,\widehat{\vect{y}}[j]\, \|\,\widehat{\vect{z}}[j] \big) \in\mathbb{F}^{3} \\ \quad o_{XYZ,j}\sample\mathbb{F} \\ \quad \cm_{XYZ,j} \gets \pedC(\ck_{XYZ},\vect{m}_{XYZ,j};o_{XYZ,j}) \\ \rt_{XYZ} \gets \merkleC((\cm_{XYZ,j})_{j\in[n]};\bot) \\
    \< \sendmessageright*[2.5cm]{\rt_{XYZ}} \< \\
    \< \< I=(j_1,\ldots,j_t)\sample[n]^t \\ \< \< \alpha_x,\alpha_y,\alpha_z\sample\mathbb{F}\setminus D \\
    \< \sendmessageleft*[2.5cm]{I, \alpha_x, \alpha_y, \alpha_z} \<  \\
    \text{For every } j\in I,\\ 
        \quad \pi_{\mathsf{mem}, j}^{ABC} \gets \merkleO(\{\cm_{ABC,\ell}\}_{\ell\in[n]},j), \\ 
        \quad \pi_{\mathsf{mem}, j}^{XYZ} \gets \merkleO(\{\cm_{XYZ,\ell}\}_{\ell\in[n]},j) \\[0.5ex]
    \pi_{\mathsf{mem}}^{ABC} \coloneqq (\cm_{ABC,j},\pi_{\mathsf{mem},j}^{ABC})_{j\in I},  \pi_{\mathsf{mem}}^{XYZ} \coloneqq (\cm_{XYZ,j},\pi_{\mathsf{mem},j}^{XYZ})_{j\in I} \\
    (\pi_{\cp},\widetilde{\cm}_{ABC},\widetilde{\cm}_{XYZ}) \gets \bPi_{\cp}.\prove(\para,\stmt,\wit) \\
    \pi_{\mathsf{link}}^{ABC} \gets \bPi_{\mathsf{link}}.\prove( \widetilde{\cm}_{ABC}, (\cm_{ABC,j})_{j\in I}; (\vect{m}_{ABC,j})_{j\in I}) \\
    \pi_{\mathsf{link}}^{XYZ} \gets \bPi_{\mathsf{link}}.\prove( \widetilde{\cm}_{XYZ}, (\cm_{XYZ,j})_{j\in I}; (\vect{m}_{XYZ,j})_{j\in I}) \\
    \pi_{\protocol}
        \coloneqq
        \big(
        \pi_{\cp},
        \widetilde{\cm}_{ABC},
        \widetilde{\cm}_{XYZ},
        \pi_{\mathsf{mem}}^{ABC},
        \pi_{\mathsf{mem}}^{XYZ},
        \pi_{\mathsf{link}}^{ABC},
        \pi_{\mathsf{link}}^{XYZ}
        \big) \\
    \< \sendmessageright*[2.5cm]{\pi_\protocol}
    }
    }
} 
\caption{The \protocol\ protocol}
\label{fig:protocol}
\end{figure*}

%%%%%%%%%%%%%

The proof consists of the SNARK proof, the sampled Merkle openings, and the
CP-Link proofs.  The sampled matrix columns and vector symbols themselves are
private SNARK witnesses; they are not public proof data.  The roots and
transcript-derived challenges are public.

We separate three kinds of work.  First, the application may need to compute or
obtain the claimed output \(\mat{C}\); this native matrix multiplication cost is
shared by all approaches and is not part of the SNARK checking cost.  Second,
the prover encodes the matrices and builds the matrix commitment root.  Third,
the online proof checks the committed multiplication statement.

\paragraph{SNARK constraints}
For each queried position \(j\), the circuit computes three length-\(k\) folds:
\[
\begin{aligned}
  &\fold(\vect{r},\widehat{\mat{A}}[*,j]),\\
  &\fold(\vect{x},\widehat{\mat{B}}[*,j]),\\
  &\fold(\vect{r},\widehat{\mat{C}}[*,j]).
\end{aligned}
\]
Thus the sampled fold constraints cost \(\mathcal{O}(tk)\).  The vector
encoding checks for
\(\widehat{\vect{x}},\widehat{\vect{y}},\widehat{\vect{z}}\) are also linear in
\(k+n\) per checked out-of-domain point for the Reed--Solomon instantiation
used in the artifact.  For constant-rate codes and fixed \(t\), the dominant
in-circuit checking cost is linear in \(k\).

\paragraph{Auxiliary work}
Merkle openings contribute \(\mathcal{O}(t\log n)\) hash work and proof
material.  CP-Link contributes backend-dependent proving and verification
costs, denoted by \(E_{\mathsf{link}}(k,t)\).  These checks are outside the
SNARK circuit.

\paragraph{Comparison}
Na\"ive SNARK arithmetization of matrix multiplication costs
\(\mathcal{O}(k^3)\) constraints.  A Freivalds-in-SNARK baseline computes the
vector-matrix products directly and costs \(\mathcal{O}(k^2)\) constraints.
\(\protocol\) replaces the full vector-matrix computation with sampled
encoded-column checks, giving \(\mathcal{O}(tk)\) fold constraints plus
auxiliary opening and linking work.

% \begin{table}[t]
% \centering
% \small
% \setlength{\tabcolsep}{3pt}
% \begin{tabularx}{\columnwidth}{@{}p{0.30\columnwidth}>{\raggedright\arraybackslash}X>{\raggedright\arraybackslash}p{0.22\columnwidth}@{}}
% \toprule
% \textbf{Component} & \textbf{Prover} & \textbf{Verifier} \\
% \midrule
% Matrix encoding/root & \(\mathcal{O}(kn)\) & roots fixed before \(r\) \\
% Intermediate vectors & \(\mathcal{O}(k^2)\) native work & --- \\
% Vector encoding/root & \(\mathcal{O}(n)\) for constant-rate codes & --- \\
% SNARK circuit & \(\mathcal{O}(tk)\) constraints & SNARK verification \\
% Merkle openings & \(\mathcal{O}(t\log n)\) & \(\mathcal{O}(t\log n)\) \\
% CP-Link & \(E_{\mathsf{link}}(k,t)\) & backend-dependent \\
% \bottomrule
% \end{tabularx}
% \caption{Asymptotic costs for \(\protocol\).  The table separates native
% matrix computation and commitment generation from the in-circuit checking
% relation.}
% \label{tab:cost-summary}
% \end{table}

\begin{table}[t]
\centering
\footnotesize
\setlength{\tabcolsep}{4pt}
\begin{tabularx}{\columnwidth}{@{}l>{\raggedright\arraybackslash}X>{\raggedright\arraybackslash}p{0.22\columnwidth}@{}}
\toprule
\textbf{Component} & \textbf{Prover} & \textbf{Verifier} \\
\midrule
Matrix root & \(\mathcal{O}(kn)\) encoding and root generation & fixed before \(r\) \\
Intermediate & \(\mathcal{O}(k^2)\) native vector products & --- \\
Vector root & \(\mathcal{O}(n)\) encoding and root generation & --- \\
CP-SNARK & \(\mathcal{O}(tk)\) constraints & proof verification \\
Merkle paths & \(\mathcal{O}(t\log n)\) & \(\mathcal{O}(t\log n)\) \\
CP-Link & \(E_{\mathsf{link}}(k,t)\) & backend dep. \\
\bottomrule
\end{tabularx}
\caption{Asymptotic costs for \(\protocol\).  Matrix and vector roots include
encoding and Merkle-root generation; CP-SNARK denotes the in-circuit sampled
checking relation.}
\label{tab:cost-summary}
\end{table}